\chapter{2012年安徽卷（理）}

\section{单选题}

\begin{problem}
复数 \(z\) 满足 \((z - \i)(2 - \i) = 5\)，则 \(z =\)（\quad）

\choices
    {\(-2 - 2\i\)}
    {\(-2 + 2\i\)}
    {\(2 - 2\i\)}
    {\(2 + 2\i\)}
\end{problem}

\begin{answer}
D
\end{answer}

\begin{solution}
由 \((z-\i)(2-\i)=5\)，得
\(z-\i=\dfrac{5}{2-\i}=\dfrac{5(2+\i)}{5}=2+\i\)，所以 \(z=2+2\i\). 故选 D.
\end{solution}

\begin{problem}
下列函数中，不满足 \( f(2x) = 2f(x) \) 的是（\quad）

\choices
    {\(f(x) = |x|\)}
    {\(f(x) = x - |x|\)}
    {\(f(x) = x + 1\)}
    {\(f(x) = -x\)}
\end{problem}

\begin{answer}
C
\end{answer}

\begin{solution}
逐项检验条件 \(f(2x)=2f(x)\). 对于 \(f(x)=|x|\)，有 \(f(2x)=|2x|=2|x|=2f(x)\)；对于 \(f(x)=x-|x|\)，有 \(f(2x)=2x-|2x|=2(x-|x|)=2f(x)\)；对于 \(f(x)=x+1\)，有 \(f(2x)=2x+1\)，而 \(2f(x)=2x+2\)，二者不相等；对于 \(f(x)=-x\)，有 \(f(2x)=-2x=2f(x)\). 因此不满足条件的是 C.
\end{solution}

\begin{problem}
如图所示，程序框图（算法流程图）的输出结果是（\quad）

\bitmapfigure[width=6.5cm,max width=0.70\linewidth]{img/2012/anhui_science/q3_fig1.png}

\choices
    {3}
    {4}
    {5}
    {8}
\end{problem}

\begin{answer}
B
\end{answer}

\begin{solution}
程序开始时 \(x=1\)，\(y=1\). 当 \(x\leqslant4\) 时执行 \(x=2x\)，\(y=y+1\)，依次得到 \((x,y)=(2,2)\)、\((4,3)\)、\((8,4)\). 此时 \(x=8>4\)，循环结束，输出 \(y=4\)，故选 B.
\end{solution}

\begin{problem}
公比为 \(\sqrt[3]{2}\) 的等比数列 \(\{a_{n}\}\) 的各项都是正数，且 \(a_{3}a_{11} = 16\)，则 \(\log_2a_{16} =\)（\quad）

\choices
    {4}
    {5}
    {6}
    {7}
\end{problem}

\begin{answer}
B
\end{answer}

\begin{solution}
设等比数列的首项为 \(a_1\)，公比为 \(q=\sqrt[3]{2}\). 由等比数列性质，
\(a_3a_{11}=a_1^2q^{12}=16a_1^2=16\). 因各项为正数，得 \(a_1=1\). 因此 \(a_{16}=q^{15}=\left(\sqrt[3]{2}\right)^{15}=2^5=32\)，所以 \(\log_2a_{16}=5\)，故选 B.
\end{solution}

\begin{problem}
甲，乙两人在一次射击比赛中各射靶 \(5\) 次，两人成绩的条形统计图如图所示，则（\quad）

\bitmapfigure[max width=0.82\linewidth]{img/2012/anhui_science/q5_fig1.png}

\choices
    {甲的成绩的平均数小于乙的成绩的平均数}
    {甲的成绩的中位数等于乙的成绩的中位数}
    {甲的成绩的方差小于乙的成绩的方差}
    {甲的成绩的极差小于乙的成绩的极差}
\end{problem}

\begin{answer}
C
\end{answer}

\begin{solution}
由甲组条形图可读出成绩为 \(4\)，\(5\)，\(6\)，\(7\)，\(8\)，故其平均数为 \(6\)，方差为
\[
\frac{(4-6)^2+(5-6)^2+(6-6)^2+(7-6)^2+(8-6)^2}{5}=2
\]
按图面柱高逐项读取，乙组成绩为
\(5\)，\(5\)，\(5\)，\(6\)，\(9\)，\(9\)，共 \(6\) 个数据，这与题干“各射靶 \(5\) 次”矛盾. 若严格按图面计算，乙组平均数为
\[
\frac{3\cdot5+6+2\cdot9}{6}=\frac{13}{2}
\]
方差为
\[
\frac{3(5-\frac{13}{2})^2+(6-\frac{13}{2})^2+2(9-\frac{13}{2})^2}{6}
=\frac{13}{4}
\]
此时甲组平均数 \(6<\frac{13}{2}\)，且甲组方差 \(2<\frac{13}{4}\)，选项 A、C 同时成立，不能得到唯一的单选结论.

若按题干规定的五次射击，并取乙组数据为
\(5\)，\(5\)，\(5\)，\(6\)，\(9\)，所以其平均数为
\[
\frac{3\cdot5+6+9}{5}=6
\]
方差为
\[
\frac{3(5-6)^2+(6-6)^2+(9-6)^2}{5}=\frac{12}{5}
\]
两组平均数相等，中位数分别为 \(6\) 和 \(5\)，极差都为 \(4\)，而
\[
2<\frac{12}{5}
\]
在这一符合题干次数的题意解释下，只有“甲的成绩的方差小于乙的成绩的方差”成立，故选 C；但图面数据与题干的矛盾须保留为版本风险.
\end{solution}

\begin{problem}
设平面 \(\alpha\) 与平面 \(\beta\) 相交于直线 \(m\)，直线 \(a\) 在平面 \(\alpha\) 内，直线 \(b\) 在平面 \(\beta\) 内，且 \(b \perp m\)，则 “\(\alpha \perp \beta\)” 是 “\(a \perp b\)” 的（\quad）

\choices
    {充分不必要条件}
    {必要不充分条件}
    {充分必要条件}
    {既不充分也不必要条件}
\end{problem}

\begin{answer}
A
\end{answer}

\begin{solution}
若 \(\alpha\perp\beta\)，则平面 \(\beta\) 内垂直于交线 \(m\) 的直线 \(b\) 垂直于平面 \(\alpha\)，从而 \(b\perp a\). 因此 \(\alpha\perp\beta\) 是 \(a\perp b\) 的充分条件.

反之，取 \(a=m\)，则由已知 \(b\perp m\) 可得 \(a\perp b\)，但 \(\alpha\) 与 \(\beta\) 不一定垂直，所以它不是必要条件. 故选 A.
\end{solution}

\begin{problem}
\((x^{2} + 2)\left(\frac{1}{x^{2}} -1\right)^{5}\) 的展开式的常数项是（\quad）

\choices
    {\(-3\)}
    {\(-2\)}
    {2}
    {3}
\end{problem}

\begin{answer}
D
\end{answer}

\begin{solution}
将 \(\left(\dfrac{1}{x^2}-1\right)^5\) 展开时，乘以 \(x^2\) 后的常数项来自 \(\left(\dfrac{1}{x^2}\right)^1\) 项，系数为 \(\mathrm C_5^1(-1)^4=5\)；乘以 \(2\) 后的常数项来自 \((-1)^5\)，系数为 \(-2\). 因此原式的常数项为 \(5-2=3\)，故选 D.
\end{solution}

\begin{problem}
在平面直角坐标系中，点 \(O(0,0)\)，\(P(6,8)\)，将向量 \(\overrightarrow{OP}\) 绕点 \(O\) 按逆时针方向旋转 \(\frac{3\pi}{4}\) 后得向量 \(\overrightarrow{OQ}\)，则点 \(Q\) 的坐标是（\quad）

\choices
    {\((-7\sqrt{2}, - \sqrt{2})\)}
    {\((-7\sqrt{2}, \sqrt{2})\)}
    {\((-4\sqrt{6}, - 2)\)}
    {\((-4\sqrt{6}, 2)\)}
\end{problem}

\begin{answer}
A
\end{answer}

\begin{solution}
旋转角为 \(\dfrac{3\pi}{4}\)，故
\(Q=\left(6\cos\dfrac{3\pi}{4}-8\sin\dfrac{3\pi}{4},6\sin\dfrac{3\pi}{4}+8\cos\dfrac{3\pi}{4}\right)=(-7\sqrt{2},-\sqrt{2})\). 故选 A.
\end{solution}

\begin{problem}
过抛物线 \( y^{2} = 4x \) 的焦点 \(F\) 的直线交该抛物线于 \(A\)，\(B\) 两点，\(O\) 为坐标原点，若 \( |AF| = 3 \)，则 \( \triangle AOB \) 的面积为（\quad）

\choices
    {\(\frac{\sqrt{2}}{2}\)}
    {\(\sqrt{2}\)}
    {\(\frac{3\sqrt{2}}{2}\)}
    {\(2 \sqrt{2}\)}
\end{problem}

\begin{answer}
C
\end{answer}

\begin{solution}
抛物线 \(y^2=4x\) 的焦点为 \(F(1,0)\). 设抛物线上参数为 \(t\) 的点为 \(M_t(t^2,2t)\). 对点 \(A=M_{t_1}\)，由焦点弦长公式 \(|AF|=1+t_1^2=3\)，得 \(t_1^2=2\). 直线 \(AB\) 过焦点时，三点 \(F\)，\(M_{t_1}\)，\(M_{t_2}\) 共线，从行列式条件得
\[
\frac{2t_1}{t_1^2-1}=\frac{2t_2}{t_2^2-1}
\]
由 \(t_1\ne t_2\)，交叉相乘得 \(t_1t_2=-1\). 由 \(t_1=\pm\sqrt{2}\)，两种取值只使图形关于 \(x\) 轴对称，三角形面积相同，故不妨取 \(t_1=\sqrt{2}\). 此时 \(t_2=-\dfrac{1}{\sqrt{2}}\)，从而 \(A=(2,2\sqrt{2})\)，\(B=\left(\dfrac{1}{2},-\sqrt{2}\right)\).

因此
\(S_{\triangle AOB}=\dfrac{1}{2}\left|x_Ay_B-y_Ax_B\right|=\dfrac{1}{2}\left|-2\sqrt{2}-\sqrt{2}\right|=\dfrac{3\sqrt{2}}{2}\)，故选 C.
\end{solution}

\begin{problem}
6位同学在毕业聚会活动中进行纪念品的交换，任意两位同学之间最多交换一次，进行交换的两位同学互赠一份纪念品. 已知6位同学之间共进行了13次交换，则收到4份纪念品的同学人数为（\quad）

\choices
    {1 或 3}
    {1 或 4}
    {2 或 3}
    {2 或 4}
\end{problem}

\begin{answer}
D
\end{answer}

\begin{solution}
把每位同学看作一个顶点，两位同学交换就在相应顶点间连一条边，则得到一个有 \(6\) 个顶点、\(13\) 条边的简单图. 每位同学收到的纪念品份数就是对应顶点的度数.

补图有 \(\mathrm C_6^2-13=15-13=2\) 条边. 在补图中，原图中度数为 \(4\) 的顶点恰好是度数为 \(1\) 的顶点. 两条边在补图中可能互不相邻，此时有 \(4\) 个度数为 \(1\) 的顶点；也可能有一个公共端点，此时有 \(2\) 个度数为 \(1\) 的顶点. 所以收到 \(4\) 份纪念品的同学人数为 \(2\) 或 \(4\)，故选 D.
\end{solution}

\section{填空题}

\begin{problem}
若 \(x\)，\(y\) 满足约束条件 \(\begin{cases}
x \geqslant 0,\\
x + 2y \geqslant 3,\\
2x + y \leqslant 3,
\end{cases}\) 则 \(x - y\) 的取值范围是\fillinblank{}.
\end{problem}

\begin{answer}
\([-3,0]\)
\end{answer}

\begin{solution}
由约束条件的边界直线得可行域的三个顶点：
\((0,\tfrac32)\)，\((0,3)\)，\((1,1)\).
目标函数为 \(x-y\)，在多边形可行域上其最值出现在顶点处. 代入三个顶点，分别得到
\(-\tfrac32\)，\(-3\)，\(0\).
因此 \(x-y\) 的取值范围是 \([-3,0]\).
\end{solution}

\begin{problem}
某几何体的三视图如图所示，该几何体的表面积是\fillinblank{}.

\bitmapfigure[max width=0.46\linewidth]{img/2012/anhui_science/q12_fig1.png}
\end{problem}

\begin{answer}
92
\end{answer}

\begin{solution}
由三视图可知，该几何体是以图示梯形为底面的直棱柱. 梯形的两条平行边长分别为 \(5\) 和 \(2\)，高为 \(4\)，斜边长为
\[
\sqrt{(5-2)^2+4^2}=5
\]
所以梯形面积为
\[
S_{\text{底}}=\frac{5+2}{2}\cdot4=14
\]
周长为 \(5+2+5+4=16\). 由三视图可得棱柱的高为 \(4\)，故表面积为
\[
2S_{\text{底}}+16\cdot4=2\cdot14+64=92
\]
\end{solution}

\begin{problem}
在极坐标系中，圆 \(\rho = 4\sin \theta\) 的圆心到直线 \(\theta = \frac{\pi}{6}\)（\(\rho \in \R\)）的距离是\fillinblank{}.
\end{problem}

\begin{answer}
\(\sqrt{3}\)
\end{answer}

\begin{solution}
将极坐标方程化为直角坐标方程：
\[
\rho=4\sin\theta
\Longrightarrow \rho^2=4\rho\sin\theta
\Longrightarrow x^2+y^2=4y
\]
因此圆的方程为 \(x^2+(y-2)^2=4\)，圆心为 \(O'(0,2)\). 直线 \(\theta=\frac{\pi}{6}\) 且 \(\rho\in\R\) 的方程为
\[
y=\frac{1}{\sqrt3}x,\qquad\text{即 }x-\sqrt3y=0
\]
圆心到该直线的距离为
\[
\frac{|0-2\sqrt3|}{\sqrt{1+3}}=\sqrt3
\]
\end{solution}

\begin{problem}
若平面向量 \(\bs{a}\)，\(\bs{b}\) 满足 \(|2\bs{a} - \bs{b}| \leqslant 3\)，则 \(\bs{a} \cdot \bs{b}\) 的最小值是\fillinblank{}.
\end{problem}

\begin{answer}
\(-\frac{9}{8}\)
\end{answer}

\begin{solution}
令 \(\bs u=2\bs a-\bs b\)，则 \(|\bs u|\leqslant3\)，且 \(\bs b=2\bs a-\bs u\). 于是
\(\bs a\cdot\bs b=2|\bs a|^2-\bs a\cdot\bs u \geqslant2|\bs a|^2-|\bs a|\,|\bs u| \geqslant2r^2-3r\)，\(r=|\bs a|\geqslant0\).
配方得
\[
2r^2-3r=2\left(r-\frac34\right)^2-\frac98\geqslant-\frac98
\]
当 \(\bs a=(\frac34,0)\)，\(\bs u=(3,0)\)，从而 \(\bs b=(-\frac32,0)\) 时，条件成立且内积等于 \(-\frac98\). 因此最小值为 \(-\frac98\).
\end{solution}

\begin{problem}
设 \(\triangle ABC\) 的内角 \(A\)，\(B\)，\(C\) 所对边的长分别为 \(a\)，\(b\)，\(c\)，则下列命题正确的是\fillinblank{}. （写出所有正确命题的编号）

\begin{circlelist}
\item 若 \(ab > c^2\)，则 \(C < \frac{\pi}{3}\).
\item 若 \(a + b > 2c\)，则 \(C < \frac{\pi}{3}\).
\item 若 \(a^3 + b^3 = c^3\)，则 \(C < \frac{\pi}{2}\).
\item 若 \((a + b)c < 2ab\)，则 \(C > \frac{\pi}{2}\).
\item 若 \((a^2 + b^2)c^2 < 2a^2b^2\)，则 \(C > \frac{\pi}{3}\).
\end{circlelist}
\end{problem}

\begin{answer}
1，2，3
\end{answer}

\begin{solution}
由余弦定理，
\[
\cos C=\frac{a^2+b^2-c^2}{2ab}
\]
逐项判断.

\begin{enumerate}
\item 若 \(ab>c^2\)，则
\[
a^2+b^2-c^2>a^2+b^2-ab\geqslant ab
\]
故 \(\cos C>\frac12\)，从而 \(C<\frac\pi3\). 命题 1 正确.

\item 若 \(a+b>2c\)，则 \(c<\frac{a+b}{2}\). 又
\[
\frac{(a+b)^2}{4}\leqslant a^2+b^2-ab
\]
因为两边之差为 \(\frac34(a-b)^2\). 因此
\[
c^2<\frac{(a+b)^2}{4}\leqslant a^2+b^2-ab
\]
故 \(\cos C>\frac12\)，从而 \(C<\frac\pi3\). 命题 2 正确.

\item 设 \(p=\frac{a^2}{a^2+b^2}\)，\(q=\frac{b^2}{a^2+b^2}\)，则 \(p+q=1\)，且 \(p\)，\(q\in(0,1)\). 所以
\[
p^{3/2}+q^{3/2}<p+q=1
\]
由 \(a^3+b^3=c^3\) 得
\[
c^3=a^3+b^3<(a^2+b^2)^{3/2}
\]
从而 \(c^2<a^2+b^2\). 于是 \(\cos C>0\)，即 \(C<\frac\pi2\). 命题 3 正确.

\item 取 \(a=b=1\)，\(c=0.9\)，三边能构成三角形，且 \((a+b)c=1.8<2=2ab\)，但
\[
\cos C=\frac{2-0.9^2}{2}>0
\]
所以 \(C<\frac\pi2\)，不是 \(C>\frac\pi2\). 命题 4 错误.

\item 取 \(a=b=1\)，\(c=\frac12\)，则 \((a^2+b^2)c^2=\frac12<2=2a^2b^2\)，但
\[
\cos C=\frac{2-\frac14}{2}>\frac12
\]
所以 \(C<\frac\pi3\)，不是 \(C>\frac\pi3\). 命题 5 错误.
\end{enumerate}

因此正确命题的编号为 1，2，3.
\end{solution}

\section{解答题}

\begin{problem}
设函数 \(f(x) = \frac{\sqrt{2}}{2}\cos \left(2x + \frac{\pi}{4}\right) + \sin^2 x\).

\begin{enumerate}
\item 求 \( f(x) \) 的最小正周期.
\item 设函数 \( g(x) \) 对任意 \( x \in \R \)，有 \( g\left(x + \frac{\pi}{2}\right) = g(x) \)，且当 \( x \in \left[0, \frac{\pi}{2}\right] \) 时，\( g(x) = \frac{1}{2} - f(x) \). 求 \( g(x) \) 在区间 \([- \pi, 0]\) 上的解析式.
\end{enumerate}
\end{problem}

\begin{solution}
\begin{enumerate}
\item 利用和差公式，
\(\frac{\sqrt2}{2}\cos\left(2x+\frac{\pi}{4}\right) =\frac12\left(\cos2x-\sin2x\right)\)，\(\sin^2x=\frac{1-\cos2x}{2}\).
因此
\[
f(x)=\frac12-\frac12\sin2x
\]
\(\sin2x\) 的最小正周期为 \(\pi\)，所以 \(f(x)\) 的最小正周期为 \(T=\pi\).
\item 由 \(g(x)=\frac12-f(x)\)，当 \(x\in[0,\frac{\pi}{2}]\) 时，
\[
g(x)=\frac12\sin2x
\]
对 \(-\pi\leqslant x<-\frac{\pi}{2}\)，有 \(x+\pi\in[0,\frac{\pi}{2})\)，故
\[
g(x)=g(x+\pi)=\frac12\sin(2x+2\pi)=\frac12\sin2x
\]
对 \(-\frac{\pi}{2}\leqslant x\leqslant0\)，有 \(x+\frac{\pi}{2}\in[0,\frac{\pi}{2}]\)，故
\[
g(x)=g\left(x+\frac{\pi}{2}\right)
=\frac12\sin\left(2x+\pi\right)
=-\frac12\sin2x
\]
所以
\[
g(x)=
\begin{cases}
\frac12\sin2x,&-\pi\leqslant x<-\frac{\pi}{2},\\
-\frac12\sin2x,&-\frac{\pi}{2}\leqslant x\leqslant0.
\end{cases}
\]
\end{enumerate}
\end{solution}

\begin{problem}
某单位招聘面试，每次从试题库中随机调用一道试题，若调用的是 A 类型试题，则使用后该试题回库，并增补一道 A 类型试题和一道 B 类型试题入库，此次调题工作结束；若调用的是 B 类型试题，则使用后该试题回库，此次调题工作结束. 试题库中现共有 \( n + m \) 道试题，其中有 \(n\) 道 A 类型试题和 \(m\) 道 B 类型试题，以 \(X\) 表示两次调题工作完成后，试题库中 A 类试题的数量.

\begin{enumerate}
\item 求 \(X = n + 2\) 的概率.
\item 设 \(m = n\)，求 \(X\) 的分布列和均值 （数学期望）.
\end{enumerate}
\end{problem}

\begin{solution}
\begin{enumerate}
\item 第一次调用 A 类试题的概率为 \(\frac{n}{n+m}\). 第一次调用 A 后，库中 A 类试题数变为 \(n+1\)，总试题数变为 \(n+m+2\)，所以第二次仍调用 A 的条件概率为 \(\frac{n+1}{n+m+2}\). 两次都调用 A 时 \(X=n+2\)，故
\[
P(X=n+2)=\frac{n}{n+m}\cdot\frac{n+1}{n+m+2}
\]
\item 当 \(m=n\) 时，第一次调用 A 或 B 的概率都为 \(\frac12\). 若第一次调用 A，第二次调用 A 的概率为
\[
\frac{n+1}{2n+2}=\frac12;
\]
若第一次调用 B，第二次调用 A 的概率为
\[
\frac{n}{2n}=\frac12
\]
因此两次调用中 A 的次数 \(X-n\) 服从参数为 \(2\)，\(\frac12\) 的二项分布，故
\[
\begin{cases}
P(X=n)=\frac14,\\
P(X=n+1)=\frac12,\\
P(X=n+2)=\frac14.
\end{cases}
\]
其均值为
\[
E(X)=n\cdot\frac14+(n+1)\cdot\frac12+(n+2)\cdot\frac14=n+1
\]
\end{enumerate}
\end{solution}

\begin{problem}
平面图形 \(ABB_{1}A_{1}C_{1}C\) 如图1所示，其中 \(BB_{1}C_{1}C\) 是矩形，\(BC = 2\)，\(BB_{1} = 4\)，\(AB = AC = \sqrt{2}\)，\(A_{1}B_{1} = A_{1}C_{1} = \sqrt{5}\). 现将该平面图形分别沿 \(BC\) 和 \(B_{1}C_{1}\) 折叠，使 \(\triangle ABC\) 与 \(\triangle A_{1}B_{1}C_{1}\) 所在平面都与平面 \(BB_{1}C_{1}C\) 垂直，再分别连接 \(A_{1}A\)，\(A_{1}B\)，\(A_{1}C_{1}\) 得到如图2所示的空间图形，对此空间图形解答下列问题.

\begin{enumerate}
\item 证明： \( AA_{1} \perp BC\).
\item 求 \(AA_{1}\) 的长.
\item 求二面角 \( A - BC - A_{1} \) 的余弦值.
\end{enumerate}

\bitmapfigure[max width=0.58\linewidth]{img/2012/anhui_science/q18_fig1.png}
\end{problem}

\begin{solution}
\begin{enumerate}
\item 建立直角坐标系，使矩形 \(BB_{1}C_{1}C\) 在 \(z=0\) 平面内，取
\(B=(0,0,0)\)，\(C=(2,0,0)\)，\(B_{1}=(0,4,0)\)，\(C_{1}=(2,4,0)\).
按图2，折叠后 \(A\) 与 \(A_{1}\) 位于平面 \(z=0\) 的两侧. 由 \(AB=AC=\sqrt2\)，且三角形 \(ABC\) 所在平面垂直于 \(z=0\) 平面，可取
\[
A=(1,0,1)
\]
同理，由 \(A_{1}B_{1}=A_{1}C_{1}=\sqrt5\)，并取与 \(A\) 异侧的折叠位置
\[
A_{1}=(1,4,-2)
\]
于是
\(\overrightarrow{AA_{1}}=(0,4,-3)\)，\(\overrightarrow{BC}=(2,0,0)\)
从而
\[
\overrightarrow{AA_{1}}\cdot\overrightarrow{BC}=0
\]
因此 \(AA_{1}\perp BC\).
\item
\[
AA_{1}=\sqrt{0^2+4^2+(-3)^2}=5
\]
\item 设 \(D\) 为 \(BC\) 的中点，则 \(D=(1,0,0)\)，且
\(\overrightarrow{DA}=(0,0,1)\)，\(\overrightarrow{DA_{1}}=(0,4,-2)\).
这两条向量分别位于二面角的两个面内，且都垂直于棱 \(BC\)，所以二面角 \(A-BC-A_{1}\) 的平面角余弦为
\[
\frac{\overrightarrow{DA}\cdot\overrightarrow{DA_{1}}}
{|\overrightarrow{DA}|\,|\overrightarrow{DA_{1}}|}
=\frac{-2}{1\cdot\sqrt{20}}
=-\frac{\sqrt5}{5}
\]
\end{enumerate}
\end{solution}

\begin{problem}
设函数 \( f(x) = a\e^{x} + \frac{1}{a\e^{x}} + b  \) （\(a > 0\)）.

\begin{enumerate}
\item 求 \( f(x) \) 在 \( [0, +\infty) \) 内的最小值.
\item 设曲线 \( y = f(x) \) 在点 \( (2, f(2)) \) 处的切线方程为 \( y = \frac{3}{2} x \)，求 \(a\)，\(b\) 的值.
\end{enumerate}
\end{problem}

\begin{solution}
\begin{enumerate}
\item 令 \(t=a\e^x\). 因为 \(a>0\) 且 \(x\geqslant0\)，所以 \(t\geqslant a\)，并且
\[
f(x)=t+\frac1t+b
\]
函数 \(t+\frac1t\) 的导数为 \(1-\frac1{t^2}\)，故它在 \((0,1]\) 上递减，在 \([1,+\infty)\) 上递增. 因此
\[
\min_{x\in[0,+\infty)}f(x)=
\begin{cases}
2+b,&0<a\leqslant1,\\
a+\dfrac1a+b,&a>1.
\end{cases}
\]
当 \(a=1\) 时两式相同，最小值在 \(x=0\) 处取得；当 \(0<a<1\) 时在 \(t=1\)，即 \(x=\ln\frac1a\) 处取得.
\item 因切线为 \(y=\frac32x\)，点 \((2,f(2))\) 在切线上，故
\[
f(2)=3
\]
又 \(f'(x)=a\e^x-\frac1{a\e^x}\)，所以
\[
f'(2)=\frac32
\]
令 \(t=a\e^2>0\)，则
\(t+\frac1t+b=3\)，\(t-\frac1t=\frac32\).
由
\[
\left(t+\frac1t\right)^2
=\left(t-\frac1t\right)^2+4
=\frac{25}{4}
\]
及 \(t+\frac1t>0\)，得 \(t+\frac1t=\frac52\). 联立
\(t+\frac1t=\frac52\)，\(t-\frac1t=\frac32\)
得到 \(t=2\)，\(\frac1t=\frac12\). 从而
\(a=\frac2{\e^2}\)，\(b=3-\frac52=\frac12\).
\end{enumerate}
\end{solution}

\begin{problem}
如图，点 \(F_{1}(-c,0)\)，\(F_{2}(c,0)\) 分别是椭圆 \(C:\frac{x^{2}}{a^{2}}+\frac{y^{2}}{b^{2}}=1\)（\(a>b>0\)）的左、右焦点，过点 \(F_{1}\) 作 \(x\) 轴的垂线交椭圆 \(C\) 的上半部分于点 \(P\)，过点 \(F_{2}\) 作直线 \(PF_{2}\) 的垂线交直线 \(x=\frac{a^{2}}{c}\) 于点 \(Q\).

\begin{enumerate}
\item 如果点 \(Q\) 的坐标是 \((4,4)\)，求此时椭圆 \(C\) 的方程.
\item 证明：直线 \(PQ\) 与椭圆 \(C\) 只有一个交点.
\end{enumerate}

\bitmapfigure[max width=0.42\linewidth]{img/2012/anhui_science/q19_fig1.png}
\end{problem}

\begin{solution}
\begin{enumerate}
\item 由于 \(F_1\) 的横坐标为 \(-c\)，直线 \(x=-c\) 与椭圆上半部分交于
\[
P=\left(-c,\frac{b^2}{a}\right)
\]
直线 \(PF_2\) 的斜率为
\[
\frac{\frac{b^2}{a}-0}{-c-c}=-\frac{b^2}{2ac}
\]
所以过 \(F_2\) 的垂线斜率为 \(\frac{2ac}{b^2}\). 与直线 \(x=\frac{a^2}{c}\) 相交时，
\[
y=\frac{2ac}{b^2}\left(\frac{a^2}{c}-c\right)
=\frac{2a(a^2-c^2)}{b^2}=2a
\]
其中用了 \(a^2=b^2+c^2\). 因而
\[
Q=\left(\frac{a^2}{c},2a\right)
\]
由 \(Q=(4,4)\) 得 \(2a=4\)，即 \(a=2\)，且 \(\frac{a^2}{c}=4\)，即 \(c=1\). 所以
\[
b^2=a^2-c^2=3
\]
椭圆方程为
\[
\frac{x^2}{4}+\frac{y^2}{3}=1
\]
\item 由
\(P=\left(-c,\frac{b^2}{a}\right)\)，\(Q=\left(\frac{a^2}{c},2a\right)\)
并利用 \(a^2=b^2+c^2\)，直线 \(PQ\) 的斜率为
\[
\frac{2a-\frac{b^2}{a}}{\frac{a^2}{c}+c}
=\frac{\frac{2a^2-b^2}{a}}{\frac{a^2+c^2}{c}}
=\frac{c}{a}
\]
故直线 \(PQ\) 的方程为
\[
y-\frac{b^2}{a}=\frac{c}{a}(x+c),
\qquad\text{即}\qquad
ay=cx+a^2
\]
将 \(y=\frac{cx+a^2}{a}\) 代入椭圆方程，得
\[
b^2x^2+(cx+a^2)^2=a^2b^2
\]
整理并使用 \(a^2=b^2+c^2\)，有
\[
a^2x^2+2a^2cx+a^4-a^2b^2
=a^2(x+c)^2=0
\]
因此 \(x=-c\)，代回直线方程得 \(y=\frac{b^2}{a}\)，交点唯一且正是 \(P\). 所以直线 \(PQ\) 与椭圆 \(C\) 只有一个交点.
\end{enumerate}
\end{solution}

\begin{problem}
数列 \(\{x_{n}\}\) 满足 \(x_{1} = 0\)，\(x_{n + 1} = -x_{n}^{2} + x_{n} + c \) （\(n \in \mathbf{N}^{*}\)）.

\begin{enumerate}
\item 证明： \(\{x_{n}\}\) 是递减数列的充分必要条件是 \(c < 0\).
\item 求 \(c\) 的取值范围，使 \( \{x_{n}\} \) 是递增数列.
\end{enumerate}
\end{problem}

\begin{solution}
\begin{enumerate}
\item 由递推式，
\[
x_{n+1}-x_n=c-x_n^2
\]
若 \(c<0\)，则对任意 \(n\)，都有 \(x_{n+1}-x_n<0\)，所以数列递减.
反之，若数列递减，则 \(x_2<x_1=0\). 由 \(x_2=c\)，可得 \(c<0\). 因此数列递减的充分必要条件是 \(c<0\).
\item 先证必要性. 若数列递增，则 \(x_2=c>x_1=0\)，故 \(c>0\). 又因所有 \(x_n\geqslant0\)，且
\[
x_{n+1}>x_n\Longrightarrow x_n^2<c
\]
所以 \(0\leqslant x_n<\sqrt c\). 数列递增且有上界，故收敛. 设其极限为 \(L\)，由递推式的连续性，
\[
L=-L^2+L+c
\]
从而 \(L=\sqrt c\).

若 \(c>\frac14\)，则 \(L>\frac12\). 因而存在 \(N\)，使 \(x_N>\frac12\). 函数
\[
h(t)=-t^2+t+c
\]
在 \((\frac12,+\infty)\) 上严格递减. 但数列递增给出 \(x_{N+1}>x_N>\frac12\)，于是
\[
x_{N+2}=h(x_{N+1})<h(x_N)=x_{N+1}
\]
矛盾. 所以 \(c\leqslant\frac14\). 必要条件为 \(0<c\leqslant\frac14\).

再证充分性. 令 \(r=\sqrt c\)，则 \(0<r\leqslant\frac12\). 用归纳法证明 \(0\leqslant x_n<r\). 初始时 \(x_1=0\). 若 \(0\leqslant x_n<r\)，则
\[
x_{n+1}-x_n=c-x_n^2>0
\]
此外，\(h(t)=-t^2+t+c\) 在 \([0,r]\) 上递增，因为 \(h'(t)=1-2t\geqslant0\)，故
\[
x_{n+1}=h(x_n)<h(r)=r
\]
所以 \(0\leqslant x_{n+1}<r\)，归纳成立，并且每一步 \(x_{n+1}>x_n\). 因此数列递增.
\[
0<c\leqslant\frac14
\]
\end{enumerate}
\end{solution}
